\documentclass[11pt, a4paper]{article}
\usepackage[a4paper, top=2.5cm, bottom=2.5cm, left=2cm, right=2cm]{geometry}
\usepackage{fontspec}
\usepackage[english, bidi=basic, provide=*]{babel}
\babelprovide[import, onchar=ids fonts]{english}
\babelfont{rm}{TeX Gyre Termes}
\babelfont{sf}{TeX Gyre Heros}
\babelfont{tt}{TeX Gyre Cursor}
\usepackage{enumitem}
\usepackage{amsmath}

\title{Comprehensive Exam Notes: Metabolism \& Cellular Respiration}
\author{Based on Chapters 8 \& 9}
\date{}

\begin{document}

\maketitle

\section*{Chapter 8: An Introduction to Metabolism}

\subsection*{8.1 Transformation of Matter and Energy}
\begin{itemize}
    \item \textbf{Metabolism:} The totality of an organism's chemical reactions; an emergent property of life.
    \item \textbf{Metabolic Pathways:} Series of steps catalyzed by specific enzymes.
    \begin{itemize}
        \item \textbf{Catabolic:} Break down complex molecules (e.g., cellular respiration); release energy ("downhill").
        \item \textbf{Anabolic:} Build complex molecules (e.g., protein synthesis); consume energy ("uphill").
    \end{itemize}
    \item \textbf{Bioenergetics:} The study of how energy flows through living organisms.
    \item \textbf{Forms of Energy:}
    \begin{itemize}
        \item \textbf{Kinetic:} Associated with motion. \textbf{Thermal energy} is kinetic energy from random atomic movement (heat in transfer).
        \item \textbf{Potential:} Based on location or structure. \textbf{Chemical energy} is potential energy available for release in a reaction.
    \end{itemize}
    \item \textbf{Thermodynamics:} 
    \begin{itemize}
        \item \textbf{First Law:} Energy is constant; can be transformed/transferred but not created/destroyed.
        \item \textbf{Second Law:} Every transfer increases \textbf{entropy} (disorder). Heat is the most disordered form of energy.
    \end{itemize}
\end{itemize}

\subsection*{8.2 Free Energy and Spontaneity}
\begin{itemize}
    \item \textbf{Gibbs Free Energy ($G$):} Portion of energy that can perform work under uniform $T$ and $P$.
    \item \textbf{Equation:} $\Delta G = \Delta H - T \Delta S$. 
    \item \textbf{Spontaneous Processes:} Occur without energy input; must have a negative $\Delta G$.
    \item \textbf{Exergonic Reactions:} Net release of free energy ($\Delta G < 0$); spontaneous.
    \item \textbf{Endergonic Reactions:} Absorbs free energy ($\Delta G > 0$); nonspontaneous.
    \item \textbf{Metabolic Equilibrium:} A cell at equilibrium is dead ($\Delta G = 0$). Cells are open systems with constant material flux.
\end{itemize}

\subsection*{8.3 ATP and Energy Coupling}
\begin{itemize}
    \item \textbf{Work Types:} Chemical (pushing reactions), Transport (pumping), Mechanical (beating cilia).
    \item \textbf{Energy Coupling:} Using exergonic processes (ATP hydrolysis) to drive endergonic ones.
    \item \textbf{ATP Hydrolysis:} $ATP + H_2O \rightarrow ADP + P_i$ ($\Delta G = -7.3$ kcal/mol).
    \item \textbf{Phosphorylated Intermediate:} Recipient of a phosphate group, more reactive/less stable.
\end{itemize}

\subsection*{8.4 Enzymes and Energy Barriers}
\begin{itemize}
    \item \textbf{Activation Energy ($E_A$):} Initial energy to start a reaction. Enzymes \textbf{lower} $E_A$ but do \textbf{not} change $\Delta G$.
    \item \textbf{Catalytic Cycle:} Substrate binds to \textbf{Active Site} $\rightarrow$ \textbf{Induced Fit} $\rightarrow$ Products released.
    \item \textbf{Factors:} Optimal temp and pH. \textbf{Cofactors} (inorganic) and \textbf{Coenzymes} (organic, vitamins).
    \item \textbf{Inhibitors:} \textbf{Competitive} (bind to active site) vs. \textbf{Noncompetitive} (bind elsewhere, change shape).
\end{itemize}

\subsection*{8.5 Regulation}
\begin{itemize}
    \item \textbf{Allosteric Regulation:} Binding at one site affects function at another.
    \item \textbf{Cooperativity:} Substrate binding at one subunit stabilizes active form for all subunits.
    \item \textbf{Feedback Inhibition:} End product inhibits an early enzyme to prevent resource waste.
\end{itemize}

\section*{Chapter 9: Cellular Respiration (Introduction)}

\subsection*{9.1 Catabolic Pathways and Redox}
\begin{itemize}
    \item \textbf{Aerobic Respiration:} Consumes $O_2$ and organic fuel.
    \item \textbf{Anaerobic Respiration/Fermentation:} Partial breakdown without $O_2$.
    \item \textbf{Redox Reactions:} Electron transfers.
    \begin{itemize}
        \item \textbf{Oxidation:} Loss of electrons ($e^-$).
        \item \textbf{Reduction:} Addition of electrons (reduces positive charge).
    \end{itemize}
    \item \textbf{Electronegativity:} $O_2$ is a powerful oxidizing agent. Moving $e^-$ closer to $O$ releases energy.
    \item \textbf{Fuel:} Molecules with many C-H bonds (glucose) have "hilltop" electrons with high potential energy.
\end{itemize}

\end{document}